%-----------------------------------------------------------------------------------------------------------------------------------------------%
%	The MIT License (MIT)
%
%	Copyright (c) 2019 Jan Küster
%
%	Permission is hereby granted, free of charge, to any person obtaining a copy
%	of this software and associated documentation files (the "Software"), to deal
%	in the Software without restriction, including without limitation the rights
%	to use, copy, modify, merge, publish, distribute, sublicense, and/or sell
%	copies of the Software, and to permit persons to whom the Software is
%	furnished to do so, subject to the following conditions:
%	
%	THE SOFTWARE IS PROVIDED "AS IS", WITHOUT WARRANTY OF ANY KIND, EXPRESS OR
%	IMPLIED, INCLUDING BUT NOT LIMITED TO THE WARRANTIES OF MERCHANTABILITY,
%	FITNESS FOR A PARTICULAR PURPOSE AND NONINFRINGEMENT. IN NO EVENT SHALL THE
%	AUTHORS OR COPYRIGHT HOLDERS BE LIABLE FOR ANY CLAIM, DAMAGES OR OTHER
%	LIABILITY, WHETHER IN AN ACTION OF CONTRACT, TORT OR OTHERWISE, ARISING FROM,
%	OUT OF OR IN CONNECTION WITH THE SOFTWARE OR THE USE OR OTHER DEALINGS IN
%	THE SOFTWARE.
%	
%
%-----------------------------------------------------------------------------------------------------------------------------------------------%


%============================================================================%
%
%	DOCUMENT DEFINITION
%
%============================================================================%

%we use article class because we want to fully customize the page and don't use a cv template
\documentclass[10pt,A4]{article}	


%----------------------------------------------------------------------------------------
%	ENCODING
%----------------------------------------------------------------------------------------

% we use utf8 since we want to build from any machine
\usepackage[utf8]{inputenc}		

%----------------------------------------------------------------------------------------
%	LOGIC
%----------------------------------------------------------------------------------------

% provides \isempty test
\usepackage{xstring, xifthen}

%----------------------------------------------------------------------------------------
%	FONT BASICS
%----------------------------------------------------------------------------------------

% some tex-live fonts - choose your own

%\usepackage[defaultsans]{droidsans}
%\usepackage[default]{comfortaa}
%\usepackage{cmbright}
\usepackage[default]{raleway}
%\usepackage{fetamont}
%\usepackage[default]{gillius}
%\usepackage[light,math]{iwona}
%\usepackage[thin]{roboto} 

% set font default
\renewcommand*\familydefault{\sfdefault} 	
\usepackage[T1]{fontenc}

% more font size definitions
\usepackage{moresize}

%----------------------------------------------------------------------------------------
%	FONT AWESOME ICONS
%---------------------------------------------------------------------------------------- 

% include the fontawesome icon set
\usepackage{fontawesome}

% use to vertically center content
% credits to: http://tex.stackexchange.com/questions/7219/how-to-vertically-center-two-images-next-to-each-other
\newcommand{\vcenteredinclude}[1]{\begingroup
\setbox0=\hbox{\includegraphics{#1}}%
\parbox{\wd0}{\box0}\endgroup}

% use to vertically center content
% credits to: http://tex.stackexchange.com/questions/7219/how-to-vertically-center-two-images-next-to-each-other
\newcommand*{\vcenteredhbox}[1]{\begingroup
\setbox0=\hbox{#1}\parbox{\wd0}{\box0}\endgroup}

% icon shortcut
\newcommand{\icon}[3] { 							
	\makebox(#2, #2){\textcolor{maincol}{\csname fa#1\endcsname}}
}	

% icon with text shortcut
\newcommand{\icontext}[4]{ 						
	\vcenteredhbox{\icon{#1}{#2}{#3}}  \hspace{2pt}  \parbox{0.9\mpwidth}{\textcolor{#4}{#3}}
}

% icon with website url
\newcommand{\iconhref}[5]{ 						
    \vcenteredhbox{\icon{#1}{#2}{#5}}  \hspace{2pt} \href{#4}{\textcolor{#5}{#3}}
}

% icon with email link
\newcommand{\iconemail}[5]{ 						
    \vcenteredhbox{\icon{#1}{#2}{#5}}  \hspace{2pt} \href{mailto:#4}{\textcolor{#5}{#3}}
}

%----------------------------------------------------------------------------------------
%	PAGE LAYOUT  DEFINITIONS
%----------------------------------------------------------------------------------------

% page outer frames (debug-only)
% \usepackage{showframe}		

% we use paracol to display breakable two columns
\usepackage{paracol}

% define page styles using geometry
\usepackage[a4paper]{geometry}

% remove all possible margins
\geometry{top=1cm, bottom=1cm, left=1cm, right=1cm}

\usepackage{fancyhdr}
\pagestyle{empty}

% space between header and content
% \setlength{\headheight}{0pt}

% indentation is zero
\setlength{\parindent}{0mm}

%----------------------------------------------------------------------------------------
%	TABLE /ARRAY DEFINITIONS
%---------------------------------------------------------------------------------------- 

% extended aligning of tabular cells
\usepackage{array}

% custom column right-align with fixed width
% use like p{size} but via x{size}
\newcolumntype{x}[1]{%
>{\raggedleft\hspace{0pt}}p{#1}}%


%----------------------------------------------------------------------------------------
%	GRAPHICS DEFINITIONS
%---------------------------------------------------------------------------------------- 

%for header image
\usepackage{graphicx}

% use this for floating figures
% \usepackage{wrapfig}
% \usepackage{float}
% \floatstyle{boxed} 
% \restylefloat{figure}

%for drawing graphics		
\usepackage{tikz}				
\usetikzlibrary{shapes, backgrounds,mindmap, trees}

%----------------------------------------------------------------------------------------
%	Color DEFINITIONS
%---------------------------------------------------------------------------------------- 
\usepackage{transparent}
\usepackage{color}

% primary color
\definecolor{maincol}{RGB}{ 45, 50, 90 }

% accent color, secondary
% \definecolor{accentcol}{RGB}{ 250, 150, 10 }

% dark color
\definecolor{darkcol}{RGB}{ 70, 70, 70 }

% light color
\definecolor{lightcol}{RGB}{245,245,245}


% Package for links, must be the last package used
\usepackage[hidelinks]{hyperref}

% returns minipage width minus two times \fboxsep
% to keep padding included in width calculations
% can also be used for other boxes / environments
\newcommand{\mpwidth}{\linewidth-\fboxsep-\fboxsep}
	


%============================================================================%
%
%	CV COMMANDS
%
%============================================================================%

%----------------------------------------------------------------------------------------
%	 CV LIST
%----------------------------------------------------------------------------------------

% renders a standard latex list but abstracts away the environment definition (begin/end)
\newcommand{\cvlist}[1] {
	\begin{itemize}{#1}\end{itemize}
}

%----------------------------------------------------------------------------------------
%	 CV TEXT
%----------------------------------------------------------------------------------------

% base class to wrap any text based stuff here. Renders like a paragraph.
% Allows complex commands to be passed, too.
% param 1: *any
\newcommand{\cvtext}[1] {
	\begin{tabular*}{1\mpwidth}{p{0.98\mpwidth}}
		\parbox{1\mpwidth}{#1}
	\end{tabular*}
}

%----------------------------------------------------------------------------------------
%	CV SECTION
%----------------------------------------------------------------------------------------

% Renders a a CV section headline with a nice underline in main color.
% param 1: section title
\newcommand{\cvsection}[1] {
	\vspace{14pt}
	\cvtext{
		\textbf{\LARGE{\textcolor{darkcol}{\uppercase{#1}}}}\\[-4pt]
		\textcolor{maincol}{ \rule{0.1\textwidth}{2pt} } \\
	}
}


%----------------------------------------------------------------------------------------
%	META SKILL
%----------------------------------------------------------------------------------------

% Renders a progress-bar to indicate a certain skill in percent.
% param 1: name of the skill / tech / etc.
% param 2: level (for example in years)
% param 3: percent, values range from 0 to 1
\newcommand{\cvskill}[3] {
	\begin{tabular*}{1\mpwidth}{p{0.72\mpwidth}  r}
 		\textcolor{black}{\textbf{#1}} & \textcolor{maincol}{#2}\\
	\end{tabular*}%
	
	\hspace{4pt}
	\begin{tikzpicture}[scale=1,rounded corners=2pt,very thin]
		\fill [lightcol] (0,0) rectangle (1\mpwidth, 0.15);
		\fill [maincol] (0,0) rectangle (#3\mpwidth, 0.15);
  	\end{tikzpicture}%
}


%----------------------------------------------------------------------------------------
%	 CV EVENT
%----------------------------------------------------------------------------------------

% Renders a table and a paragraph (cvtext) wrapped in a parbox (to ensure minimum content
% is glued together when a pagebreak appears).
% Additional Information can be passed in text or list form (or other environments).
% the work you did
% param 1: time-frame i.e. Sep 14 - Jan 15 etc.
% param 2:	 event name (job position etc.)
% param 3: Customer, Employer, Industry
% param 4: Short description
% param 5: work done (optional)
% param 6: technologies include (optional)
% param 7: achievements (optional)
\newcommand{\cvevent}[7] {
	
	% we wrap this part in a parbox, so title and description are not separated on a pagebreak
	% if you need more control on page breaks, remove the parbox
	\parbox{\mpwidth}{
		\begin{tabular*}{1\mpwidth}{p{0.72\mpwidth}  r}
	 		\textcolor{black}{\textbf{#2}} & \colorbox{maincol}{\makebox[0.25\mpwidth]{\textcolor{white}{#1}}} \\
			\textcolor{maincol}{\textbf{#3}} & \\
		\end{tabular*}\\[4pt]
	
		\ifthenelse{\isempty{#4}}{}{
			\cvtext{#4}\\
		}
	}

	\ifthenelse{\isempty{#5}}{}{
		\vspace{4pt}
		{#5}
	}
	\vspace{4pt}
}


%----------------------------------------------------------------------------------------
%	 CV Pub
%----------------------------------------------------------------------------------------

% Define an environment for cvpub
\newenvironment{cvpubs}{%
  \vspace{\acvSectionContentTopSkip}
    \vspace{-2mm}
    \paragraphstyle
}{%
    \par
    \vspace{1mm}
}

% Define a line of cvpub
% Usage: \cvpub{<ciation>}
\newcommand*{\cvpub}[1]{%
	\leftskip 0.7cm
    \parindent -0.7cm
    \pubstyle{#1}
}

% Define an environment for cvitems(for cventry)
\newenvironment{cvitems}{%
  \vspace{-4.0mm}
  \begin{justify}
  \begin{itemize}[leftmargin=2ex, nosep, noitemsep]
    \setlength{\parskip}{0pt}
    \renewcommand{\labelitemi}{\bullet}
}{%
  \end{itemize}
  \end{justify}
  \vspace{-4.0mm}
}

%----------------------------------------------------------------------------------------
%	 CV META EVENT
%----------------------------------------------------------------------------------------

% Renders a CV event on the sidebar
% param 1: title
% param 2: subtitle (optional)
% param 3: customer, employer, etc,. (optional)
% param 4: info text (optional)
\newcommand{\cvmetaevent}[4] {
	\textcolor{maincol} {\cvtext{\textbf{\begin{flushleft}#1\end{flushleft}}}}

	\ifthenelse{\isempty{#2}}{}{
	\textcolor{darkcol} {\cvtext{\textbf{#2}} }
	}

	\ifthenelse{\isempty{#3}}{}{
		\cvtext{{ \textcolor{darkcol} {#3} }}\\
	}

	\cvtext{#4}\\[14pt]
}

%---------------------------------------------------------------------------------------
%	QR CODE
%----------------------------------------------------------------------------------------

% Renders a qrcode image (centered, relative to the parentwidth)
% param 1: percent width, from 0 to 1
\newcommand{\cvqrcode}[1] {
	\begin{center}
		\includegraphics[width={#1}\mpwidth]{qrcode}
	\end{center}
}

%=+=+=+=+=+=+=+=+=+=+=+=+=+=+=+=+=+=+=+=+=+=+=+=+=+=+=+=+=+=+=+=+=+=+=+=+=+=+=+=+
%,,,,,,,,,,,,,,,,,,,,,,,,,,,,,,,,,,,,,,,,,,,,,,,,,,,,,,,,,,,,,,,,,,,,,,,,,,,,,,,,
                       % EDIT AFTER THIS POINT
%''''''''''''''''''''''''''''''''''''''''''''''''''''''''''''''''''''''''''''''''
%=+=+=+=+=+=+=+=+=+=+=+=+=+=+=+=+=+=+=+=+=+=+=+=+=+=+=+=+=+=+=+=+=+=+=+=+=+=+=+=+


%============================================================================%
%
%
%
%	DOCUMENT CONTENT
%
%
%
%============================================================================%
\begin{document}
\columnratio{0.31}
\setlength{\columnsep}{2.2em}
\setlength{\columnseprule}{4pt}
\colseprulecolor{lightcol}
\begin{paracol}{2}
\begin{leftcolumn}
%---------------------------------------------------------------------------------------
%	META IMAGE
%----------------------------------------------------------------------------------------
%\includegraphics[width=\linewidth]{untitled.jpg}	%trimming relative to image size


\vfill\null
\cvsection{CONTACT}
	
\iconemail{EnvelopeSquare}{14}{ecnucrystal@gmail.com}{ecnucrystal@gmail.com}{black}\\[6pt]
\icontext{Github}{14}{\href{https://github.com/Jing-L97}{@Jing-L97}}{black}\\[6pt]
\icontext{Linkedin}{14}{\href{https://www.linkedin.com/in/jing-liu-compuling/}{@Jing Liu}}{black}\\[6pt]
\icontext{Phone}{14}{+33 758 828 907}{black}\\[6pt]
\vfill\null
%\cvqrcode{0.7}


%\vfill\null
%\cvqrcode{0.7}



%---------------------------------------------------------------------------------------
%	META SKILLS
%----------------------------------------------------------------------------------------
\cvsection{SUMMARY}
\vspace{-1cm}
\cvmetaevent
{}
{}
% {PhD researcher (Marie Curie Fellow) working on cognitively grounded and data-efficient language modeling. Specialized in LLM \textbf{mechanistic interpretability}, \textbf{retrieval-augmented generation}, and \textbf{conversational agents}. Published at NeurIPS, ICML, CogSci, and ICMI. Co-organizer of \textbf{BabyLM 2025}, \textbf{BabyLM-Speech 2026}.}
{PhD researcher and \textbf{Marie Skłodowska-Curie Fellow} working on large language models: \textbf{mechanistic interpretability}, \textbf{RL fine-tuning} (PPO/DPO/GRPO), \textbf{retrieval augmentation}, and \textbf{automatic evaluation}. I build data-efficient training and evaluation pipelines and turn them into reusable benchmarks and baselines adopted by the community (BabyLM 2025). First-author work at NeurIPS, ICML, COLM, and CogSci. Area Chair at EMNLP; co-organizer of BabyLM 2025/2026.}
{}

%---------------------------------------------------------------------------------------
%	META SKILLS
%----------------------------------------------------------------------------------------

\cvsection{SKILLS}
\vspace{6pt}

\textbf{\textcolor{darkcol}{Languages}}\\[2pt]
\cvtext{Python, R, MATLAB, Java, Prolog, JavaScript, HTML, Bash}\\[6pt]

\textbf{\textcolor{darkcol}{ML \& DL Frameworks}}\\[2pt]
\cvtext{PyTorch, TensorFlow, Fairseq, Hugging Face TRL, vLLM}\\[6pt]

\textbf{\textcolor{darkcol}{LLMs \& Alignment}}\\[2pt]
\cvtext{RLHF (PPO, DPO, GRPO), reward modeling, RAG, LLM-as-a-judge evaluation}\\[6pt]

\textbf{\textcolor{darkcol}{Interpretability}}\\[2pt]
\cvtext{Logit attribution (LogitLens), circuit tracing, sparse autoencoders (SAE)}\\[6pt]

\textbf{\textcolor{darkcol}{Speech \& Audio}}\\[2pt]
\cvtext{Self-supervised speech models (HuBERT, wav2vec 2.0), speaker diarization (pyannote), forced alignment, prosody (openSMILE), Praat, ELAN}\\[6pt]

\textbf{\textcolor{darkcol}{Methods}}\\[2pt]
\cvtext{Benchmark \& evaluation design, ablation studies, statistical analysis, data-efficient / curriculum training}\\[6pt]

\textbf{\textcolor{darkcol}{Infra \& Tooling}}\\[2pt]
\cvtext{Git, Linux, Weights \& Biases, LaTeX}

\vspace{1.8cm}

%---------------------------------------------------------------------------------------
%	ACHIEVEMENTS
%----------------------------------------------------------------------------------------
% \newpage
\cvsection{ACHIEVEMENTS}
\vspace{6pt}

\textbf{\textcolor{darkcol}{Fellowships}}\\[5pt]
Marie-Curie PhD Fellow (2023)\\[4pt]
Radboud Honors Academy\\[4pt]
National Research Training Program

\vspace{8pt}
\textbf{\textcolor{darkcol}{Scholarships \& Awards}}\\[5pt]
NeurIPS Travel Fund (2025)\\[4pt]
INTERSPEECH Student Grant\\[4pt]
Radboud Scholarship\\[4pt]
2nd Prize – National Teaching Competition\\[4pt]
ECNU Guanghua Scholarship\\[4pt]
ECNU Scholarship


\end{leftcolumn}
\begin{rightcolumn}
%---------------------------------------------------------------------------------------
%	TITLE  HEADER
%----------------------------------------------------------------------------------------
\fcolorbox{white}{darkcol}{\begin{minipage}[c][3.5cm][c]{1\mpwidth}
	\begin {center}
		\HUGE{ \textbf{ \textcolor{white}{ \uppercase{ Jing LIU } } } } \\[-24pt]
		\textcolor{white}{ \rule{0.1\textwidth}{1.25pt} } \\[4pt]
		% \large{ \textcolor{white} {PhD Researcher - AI + Cognitive Science } }
		\large{ \textcolor{white} {PhD Researcher --- LLMs $\cdot$ Interpretability $\cdot$ RL \& Evaluation } }
	\end {center}
\end{minipage}} \\[14pt]
\vspace{-12pt}

%---------------------------------------------------------------------------------------
%	EDUCATION
%----------------------------------------------------------------------------------------
%\vfill\null
\cvsection{EDUCATION}

\cvevent
	{\textbf{2023 - ongoing}}
	{Ph. D. -  NLP \& Cognitive Modeling}
	{École Normale Supérieure, Paris, France}
	{Thesis: Data-Efficient Spoken Language Modeling (Advisors: Emmanuel Dupoux, Najoung Kim, David Hawarth)}
\vfill\null

\cvevent
	{\textbf{2022 - 2023}}
	{Adv. Master - Artificial Intelligence}
	{KU Leuven, Leuven (Belgium)}
	{Thesis: Computational Modeling of Early Word Acquisition.}
\vfill\null

\cvevent
	{\textbf{2020 - 2022}}
	{Research Master - Linguistics \& Comm. Sciences (Distinction)}
	{Radboud University, Nijmegen (the Netherlands)}
	{Thesis: Predicting Backchannels in Multimodal Child-Caregiver Conversation (Advisors: Abdellah Fourtassi, Asli Özyürek) }
\vfill\null


\cvevent
	{\textbf{2015 - 2019}}
	{BA in English education \& BSc in Psychology}
	{East China Normal University, Shanghai, China}
	{Thesis: Mnemonic Effects in Idiom Learning. (Advisor: Yanning Yang)}
\vfill\null


%---------------------------------------------------------------------------------------
%	WORK
%----------------------------------------------------------------------------------------
%\vfill\null
\cvsection{WORK}

\cvevent
	{\textbf{2026.7 - 2026.12}}
	{applied scientist intern}
	{Amazon, Berlin, Germany}
	{Anomaly detection}
\vfill\null

\cvevent
	{\textbf{2022.2 - 2022.8}}
	{research intern}
	{Inria, Paris, France}
	{Conversation behavior prediction}
\vfill\null

\cvevent
	{\textbf{2019 - 2020}}
	{English teacher in senior high school}
	{Hohhot, Inner Mongolia, China}
	{English courses for senior high school students}
\vfill\null

%---------------------------------------------------------------------------------------
%	PUBLICATION
%----------------------------------------------------------------------------------------

\vfill\null

\cvsection{PUBLICATIONS}

\textbf{\textcolor{darkcol}{Peer-Reviewed (Conferences \& Workshops)}}
\begin{itemize}
	\item \textbf{Liu, J.}, \& Shukla, A. (2026). Compression at a Cost: Interpreting Information Bottlenecks in Safety-Aligned Reward Models. Mechanistic Interpretability Workshop, ICML, Seoul, South Korea.
	\item \textbf{Liu, J.}, \& Kim, N. (2026). Does Episodic Memory Help Close the Lexical Frequency Gap in Sensitivity to Syntactic Contrasts? A Test Using Retrieval-Augmented Language Models. COLM, San Francisco, USA.
	\item \textbf{Liu, J.}, Li, Y., \& Wang, H. (2025). Distributed Specialization: Rare-Token Neurons in Large Language Models. Mechanistic Interpretability Workshop, NeurIPS, San Diego, USA.
	\item \textbf{Liu, J.} (2025). No Clustering, No Routing: How Transformers Actually Process Rare Tokens. UniReps Workshop, NeurIPS, San Diego, USA.
	\item \textbf{Liu, J.}, Wang, H., \& Li, Y. (2025). Emergent Specialization: Rare Token Neurons in Language Models. High-dimensional Learning Dynamics Workshop, ICML, Vancouver, Canada.
	\item \textbf{Liu, J.}, \& Fourtassi, A. (2025). Benchmarking LLMs for Mimicking Child-Caregiver Language in Interaction. Cognitive Science Society, San Francisco, USA.
	\item Charpentier, L., Choshen, L., Cotterell, R., Gul, M. O., Hu, M., Jumelet, J., Linzen, T., \textbf{Liu, J.}, Mueller, A., Ross, C., Shah, R. S., Warstadt, A., Wilcox, E., \& Williams, A. (2025). Findings of the Third BabyLM Challenge. Proceedings of the Third BabyLM Workshop, ACL.
	\item Agrawal, A., \textbf{Liu, J.}, Bodur, K., Favre, B., \& Fourtassi, A. (2023). Development of Multimodal Turn Coordination in Conversations: Evidence for Adult-like Behavior in Middle Childhood. Cognitive Science Society, Sydney, Australia.
	\item \textbf{Liu, J.}, Nikolaus, M., Bodur, K., \& Fourtassi, A. (2022). Predicting Backchannel Signaling in Multimodal Child-Caregiver Conversation. WoCBU Workshop, ICMI 2022, Bangalore, India.
\end{itemize}

\vspace{4pt}
\textbf{\textcolor{darkcol}{Preprints \& Under Review}}
\begin{itemize}
	\item \textbf{Liu, J.}, Schweitzer, S., \& Fourtassi, A. (2026). Which Forms of Caregiver Feedback Support Grammar Learning? A Reinforcement-Learning Study of Child-Like Language Models. arXiv.
	\item \textbf{Liu, J.}, \& Xie, S. (2026). Is Curriculum Learning Based on Pedagogical Child-Directed Speech Effective for Language Model Training? arXiv.
	\item \textbf{Liu, J.}, \& Dupoux, E. (2026). On Comparing Developmental Curves in Children and AI Simulations. Annual Review of Developmental Psychology (under review).
	\item \textbf{Liu, J.} (2025). Layer Specialization Underlying Compositional Reasoning in Transformers. arXiv.
	\item \textbf{Liu, J.}, Lavechin, M., Hamilakis, N., \& Dupoux, E. (2025). CORPUS-CDI: Comparing Children and AI Models' Learning Rate of Expressive Vocabulary. arXiv.
	\item \textbf{Liu, J.} (2025). Circuit-Aware Reward Training: A Mechanistic Framework for Longtail Robustness in RLHF. arXiv.
	\item \textbf{Liu, J.}, Lavechin, M., Hamilakis, N., \& Dupoux, E. (2025). Machine CDI: A Lexical Benchmark for Expressive Vocabulary (submitted).
	\item Charpentier, L., Choshen, L., Cotterell, R., Gul, M. O., Hu, M., Jumelet, J., Linzen, T., \textbf{Liu, J.}, Mueller, A., Ross, C., Shah, R. S., Warstadt, A., Wilcox, E., \& Williams, A. (2025). BabyLM Turns 3: Call for Papers for the 2025 BabyLM Workshop. arXiv.
\end{itemize}



%---------------------------------------------------------------------------------------
%	PROJECTS
%----------------------------------------------------------------------------------------
\vfill\null
% \cvsection{SELECTED PROJECTS}

% \cvevent
% 	{\textbf{2026}}
% 	{Multi-communicative feedback in child-caregiver interactions}
% 	{Ongoing}
% 	{Implemented reward aggregation and automatic weighting to align LM generations with caregiver feedback.}
% \vfill\null


% \cvevent
% 	{\textbf{2026}}
% 	{Evaluating speech language model in lexical discrimination tasks}
% 	{Ongoing}
% 	{Self supervised speech representation learning}
% \vfill\null

% \cvevent
% 	{\textbf{2025-26}}
% 	{MI-Reward: Mitigating Reward Hacking via Sparse Autoencoders}
% 	{Ongoing}
% 	{Identify and suppress spurious features in reward modeling using Sparse Autoencoders}
% \vfill\null



% \cvevent
% 	{\textbf{2025}}
% 	{Grokked transformer for in-context learning}
% 	{Ongoing}
% 	{Analyzed compositional reasoning with circuit tracing and quantum-inspired parameter efficient attention mechanisms.}
% \vfill\null


% \cvevent
% 	{\textbf{2025}}
% 	{Emergent Specialization in LLMs}
% 	{NeurIPS MechInterp Workshop, 2025; ICML HiLD Workshop, 2025}
% 	{Analyzed rare token neuron specialization using logit attribution and sparse autoencoders.}
% \vfill\null


% \cvevent
% 	{\textbf{2024-25}}
% 	{Retrieval-Augmented LMs for rare token acquisition }
% 	{Cognitive Science 2026}
% 	{Investigated factors contributing to RAG system retrieval. }
% \vfill\null


% \cvevent
% 	{\textbf{2024}}
% 	{Communicative feedback in child-caregiver interactions}
% 	{Baseline model of BabyLM Challenge 2025}
% 	{Implemented PPO-based finetuning to align LM generations with caregiver feedback. Built baseline for BabyLM Challenge 2025.}
% \vfill\null


% \cvevent
% 	{\textbf{2024}}
% 	{Benchmarking LLMs for Mimicking Child-Caregiver Language in Interaction}
% 	{Cognitive Science 2025; metrics of BabyLM Challenge 2025}
% 	{Built LLM-based conversational agents to simulate child-adult interactions. Evaluated generated interactions with multi-level linguistic properties}
% \vfill\null

% \cvevent
% 	{\textbf{2023-2024}}
% 	{Machine CDI Benchmark}
% 	{Under Review}
% 	{Built a lexical benchmark evaluating expressive vocabulary generation in LLMs vs. children. Grounded evaluation in developmental norms.}
% \vfill\null


% \cvevent
% 	{\textbf{2022-2023}}
% 	{Automatic predicting conversational behaviors}
% 	{ICMI 2022, Cognitive Science 2023}
% 	{Built multi-modal model to predict conversational behaviors like backchanneling and turn-taking}
% \vfill\null



\cvsection{SELECTED PROJECTS}
\vspace{4pt}

\textbf{\textcolor{darkcol}{LLM Alignment \& RL Fine-Tuning}}
\vspace{2pt}

\cvevent
	{\textbf{2024}}
	{Communicative-Feedback Alignment (BabyLM 2025 baseline)}
	{Official baseline, BabyLM Challenge 2025}
	{Implemented PPO-based RLHF to align child-like LM generations with caregiver feedback signals; released as an official baseline adopted by BabyLM 2025 participants.}
\vfill\null

\cvevent
	{\textbf{2026}}
	{Multi-Reward Feedback Alignment}
	{Submitted to EMNLP 2026}
	{Built reward aggregation with automatic weighting to align LM generations across multiple caregiver-feedback signals.}
\vfill\null

\vspace{4pt}
\textbf{\textcolor{darkcol}{Mechanistic Interpretability}}
\vspace{2pt}

\cvevent
	{\textbf{2025-26}}
	{MI-Reward: Mitigating Reward Hacking via Sparse Autoencoders}
	{Mechanistic interpretability Work shop ICML 2026}
	{Built an SAE-based reward-modeling pipeline to identify and suppress spurious features that drive reward hacking in RLHF.}
\vfill\null

\cvevent
	{\textbf{2025}}
	{Emergent Specialization in LLMs}
	{NeurIPS MI Workshop; ICML HiLD Workshop, 2025}
	{Characterized rare-token neuron specialization in transformers via logit attribution and sparse autoencoders, showing distributed (not modular) processing of rare tokens.}
\vfill\null

\cvevent
	{\textbf{2025}}
	{Compositional Reasoning in Grokked Transformers}
	{Ongoing}
	{Analyzed compositional in-context reasoning using circuit tracing and a parameter-efficient attention design.}
\vfill\null

\vspace{4pt}
\textbf{\textcolor{darkcol}{Evaluation \& Benchmarks}}
\vspace{2pt}

\cvevent
	{\textbf{2023-24}}
	{Machine CDI: Lexical Benchmark for LLMs}
	{Submitted}
	{Constructed a lexical benchmark grounding LLM expressive-vocabulary evaluation in developmental (CDI) norms, enabling direct child-vs-model comparison.}
\vfill\null

\cvevent
	{\textbf{2024}}
	{LLM Conversational Agents \& Evaluation Suite}
	{CogSci 2025; metrics, BabyLM 2025}
	{Built LLM-based agents simulating child--adult dialogue and a multi-level linguistic evaluation suite; contributed evaluation metrics to the BabyLM 2025 Challenge.}
\vfill\null

\cvevent
	{\textbf{2024-25}}
	{Retrieval-Augmented LMs for Rare-Token Acquisition}
	{COLM 2026}
	{Studied which retrieval factors let RAG close the lexical-frequency gap in syntactic sensitivity, using controlled retrieval-augmented evaluation.}
\vfill\null

\vspace{4pt}
\textbf{\textcolor{darkcol}{Speech \& Multimodal Modeling}}
\vspace{2pt}

\cvevent
	{\textbf{2026}}
	{Speech LM Evaluation for Lexical Discrimination}
	{Ongoing}
	{Evaluated self-supervised speech representations on lexical-discrimination tasks for spoken language models.}
\vfill\null

\cvevent
	{\textbf{2022-23}}
	{Multimodal Conversational-Behavior Prediction}
	{ICMI 2022; CogSci 2023}
	{Built multimodal models to predict backchanneling and turn-taking from audio, video, and transcripts.}
\vfill\null


%---------------------------------------------------------------------------------------
%	PERSONAL DETAILS
%----------------------------------------------------------------------------------------
\vfill\null

\cvsection{SERVICE \& LEADERSHIP}
\vspace{-0.3cm}
\begin{itemize}
  \item \textbf{Area Chair}, ACL Rolling Review (EMNLP), 2026
  \item \textbf{Co-organizer}, BabyLM Challenge 2025; BabyLM-Speech 2026; Computational Developmental Linguistics, ACL 2026
  \item \textbf{Invited Lecturer}, Boston University, Spring 2024
  \item \textbf{Reviewer}: ICML \& NeurIPS Mechanistic Interpretability / UniReps Workshops (2025--26), ACL 2025, BabyLM (ACL/EMNLP), \textit{Open Mind} journal
  \item \textbf{Organizer}, Weekly NLP Reading Group, 2024--present
  \item \textbf{Teaching}: Psycholinguistics \& Comprehensive English (2019--20); TA, English Phonetics (2017--18)
\end{itemize}

\vfill\null

% hotfixes to create fake-space to ensure the whole height is used
\vfill
\vfill
\vfill
\end{rightcolumn}
\end{paracol}
\end{document}

